\documentclass{article}
\usepackage[utf8]{inputenc}
\usepackage{amsmath}
\usepackage{geometry}
\geometry{a4paper, margin=1in}

\title{Labwork 1: Gradient Descend}
\author{Le Nguyen Minh Quang} 
\date{\today}

\begin{document}

\maketitle

\section{Algorithm Implementation}
To implement the Gradient Descent algorithm from scratch, I first defined the objective function $f(x) = x^2$ and its first-order derivative $f'(x) = 2x$ using Python functions. 

The core logic is implemented in the \texttt{gradient\_descent} function. The algorithm runs for a fixed number of iterations. In each iteration, it updates the current value of $x$ using the formula:
\begin{equation}
    x_{new} = x_{old} - r \cdot f'(x_{old})
\end{equation}
where $r$ is the learning rate. To track the progress, I print the intermediate steps including the iteration number, the current value of $x$, and the function value $f(x)$.

\section{Effect of Different Learning Rates ($r$)}
I experimented with various learning rates starting from an initial point $x = 10$. The learning rate $r$ significantly affects the convergence of the algorithm:

\begin{itemize}
    \item \textbf{Small learning rate (e.g., $r = 0.01$):} The updates to $x$ are very tiny. The algorithm is stable but converges very slowly, requiring a large number of iterations to reach the minimum ($x = 0$).
    \item \textbf{Good learning rate (e.g., $r = 0.1$):} The algorithm converges smoothly and efficiently to the minimum within a reasonable amount of iterations.
    \item \textbf{Large learning rate (e.g., $r = 1.1$):} The step size is too large. Instead of moving towards the minimum, the value of $x$ overshoots the minimum and diverges to infinity.
\end{itemize}

In conclusion, choosing an appropriate learning rate is crucial. A value too small leads to slow training, while a value too large causes the algorithm to diverge.

\end{document}
